\section{Threat Model and Attack Space}
\label{sec:threat}

This section defines the adversary, the parameters under its control, the conditions its
transmission must satisfy to capture a receiver, and the knowledge it is assumed to
possess. Every adversary action in this study is a choice of transmitter settings, so the
threat model is stated in terms of a transmitter and not in terms of a displacement in
a feature space.

\subsection{Adversary model and objective}

The adversary is a spoofer: a transmitter emitting counterfeit GPS L1 C/A signals towards
a victim receiver. Its objective is to cause the victim to report a position or a time of
the adversary's choosing while the victim continues to operate nominally, producing a
converged navigation solution with residuals of ordinary magnitude.

The adversary's only channel to the victim is the radio interface. It does not modify the
victim's software, read its internal state, or alter the observables the victim computes.
Its influence on those observables is mediated entirely by the receiver's own acquisition
and tracking of the composite signal, which is what determines the adversary's variables
in this study.

\subsection{Transmitter controlled parameters}

Table~\ref{tab:attacker} lists the parameters the adversary sets. A choice of all five
defines one attack.

\begin{table}[t]
\centering
\caption{The attacker's controls. Each row is a transmitter setting, not a property of
the victim's measurements. One choice of all five defines one attack.}
\label{tab:attacker}
\begin{tabular}{@{}llp{8.5cm}@{}}
\toprule
Setting & Symbol & What it does \\
\midrule
Power advantage      & $\alpha$   & Strength of the counterfeit relative to the
                                    authentic signal, in decibels \\
Code offset          & $\tau_0$   & Alignment of the counterfeit code at the start,
                                    in chips \\
Code pull-off rate   & $\dot\tau$ & Rate the counterfeit code is dragged away from the
                                    authentic code \\
Carrier offset       & $\Delta f$ & Frequency of the counterfeit carrier relative to
                                    the authentic carrier \\
Carrier phase        & $\phi_0$   & Phase of the counterfeit carrier relative to the
                                    authentic carrier \\
\bottomrule
\end{tabular}
\end{table}

Two of these parameters carry the attack. The power advantage $\alpha$ determines whether
the victim's tracking loops transfer to the counterfeit at all, and the pull-off rate
$\dot\tau$ determines how rapidly the victim's reported range moves once transfer has
occurred. One chip of GPS L1 C/A code corresponds to
$c/f_\mathrm{chip} = 293.05$~m of range, where $c$ is the speed of light and
$f_\mathrm{chip} = 1.023$~MHz is the chipping rate, so a pull-off rate expressed in chips
per second converts directly to metres per second of induced range error. The remaining
three parameters govern the alignment of the counterfeit at onset, and therefore govern
whether the capture conditions of Section~\ref{sec:takeover} are met.

\subsection{Receiver takeover conditions}
\label{sec:takeover}

The adversary cannot select these parameters arbitrarily and still capture a receiver that
is already tracking the authentic signals. Two conditions govern capture
\citep{psiaki2016gnss,kerns2014uav}.

The counterfeit must arrive aligned with the authentic signal in code phase, carrier phase
and Doppler frequency, so that the victim's correlator observes a single undistorted peak
and not two competing peaks. A misaligned counterfeit does not transfer the tracking
loops; it perturbs the victim's observables without displacing its position or time
solution, and therefore constitutes interference and not spoofing. The counterfeit must
additionally exceed the authentic signal in received power, since code tracking follows the
dominant correlation peak. In practice the power advantage employed is small, typically a
few decibels, because a large advantage is itself observable in the receiver's automatic
gain control and in the reported $C/N_0$. This second condition is the established one
and is adopted here, though it may not be necessary: a spoofing method has since been
reported that uses controlled destructive multipath to pull the code loop at a power
below the authentic signal \citep{liu2026dimas}. Requiring the power advantage therefore
makes the attack library used here a conservative subset of what an adversary may
transmit.

Only once both conditions hold does the adversary begin to displace the code. The
experiments in this paper begin from that point: the generator is seeded from the victim's
own tracking so that the alignment condition is satisfied by construction, and what is
measured is detector behaviour after capture has been achieved and not the difficulty
of achieving it. These
conditions define the subset of the parameter space in which an attack is operative, and
Section~\ref{sec:method} applies them to partition the swept library into transmissions
that capture the receiver and transmissions that do not. Detector performance is reported
on the former, so that no detector is credited for flagging a transmission that would not
have succeeded.

\subsection{Knowledge and access assumptions}

The adversary is assumed to have no knowledge of the victim's detector. It does not know
which of the fourteen detectors is deployed, what data that detector was trained on, what
decision threshold it uses, or what score it assigns to any observation. No gradient of any
detector is computed and no detector is queried at any point in the attack construction.

In the taxonomy of Section~\ref{sec:rw_adversarial} this is weaker than a decision based
adversary, which at least observes the classifier's output label. It corresponds to a
spoofer transmitting without feedback from the victim, which is the operationally
realistic case for an attack on a receiver whose alarm state is not externally visible. The
detection failures reported in Section~\ref{sec:results} are therefore failures against an
adversary with no information about the system it defeats.

\subsection{Scope of the adversary model}

Four further restrictions delimit the adversary modelled here relative to the capability a
spoofer may possess in general.

\begin{itemize}
  \item \textbf{No detector knowledge.} As stated above, the transmission is not shaped
  against any specific model.
  \item \textbf{All tracked satellites are spoofed.} An adversary that leaves one
  satellite authentic gains an additional degree of freedom, at the cost of an
  inconsistency between that satellite and the remainder which a cross satellite check may
  identify.
  \item \textbf{A single signal is transmitted}, GPS L1 C/A. A receiver combining multiple
  constellations or multiple frequencies retains evidence that this adversary does not
  attempt to defeat.
  \item \textbf{The victim is stationary and its position is surveyed}, so a common
  pull-off applied to all channels is absorbed into the receiver clock estimate and
  manifests as a timing error.
\end{itemize}

These restrictions determine how the results should be interpreted. Each detection failure
reported in Section~\ref{sec:results} is obtained against this restricted adversary, so the
measured vulnerability is a lower bound on the vulnerability to an adversary that relaxes
any of them. The fourth restriction additionally fixes the observable consequence measured
in Section~\ref{sec:res_nav}: because a common pull-off on all channels is absorbed by the
clock, the induced error appears in the timing solution, and it is that error which is
compared against the commanded rate.
